\chapter{Analiza problemei}
\label{ch:analiza}

\section{Definirea problemei}

Se consideră un magazin online cu \textbf{200 de categorii}, \textbf{100.000 de produse} și
\textbf{500.000 de recenzii}, stocate în PostgreSQL. Aplicația web expune trei operații:

\begin{enumerate}
  \item \texttt{GET /api/products/\{id\}} — detaliile unui produs împreună cu media notelor și numărul de
        recenzii (în forma normalizată: \texttt{LEFT JOIN} + agregare pe $\approx 5$ rânduri);
  \item \texttt{GET /api/categories/\{id\}/top} — cele mai bine cotate produse dintr-o categorie
        (în forma normalizată: \texttt{JOIN} + \texttt{GROUP BY} + sortare peste $\approx 2.500$ de rânduri);
  \item \texttt{PUT /api/products/\{id\}} — modificarea prețului și/sau stocului.
\end{enumerate}

Traficul este dominat de citiri (90–99\%), cu popularitate Zipf. Problema este: \emph{cum poate fi
servit acest trafic cu debit maxim și latență minimă, păstrând datele corecte, folosind memorie
limitată și mai multe mașini?}

\section{Cerințe}

\subsection{Cerințe funcționale}
\begin{itemize}
  \item F1. Citirea produselor și a clasamentelor, cu indicarea sursei (\texttt{X-Cache: HIT/MISS}).
  \item F2. Modificarea produselor, cu propagarea corectă în cache conform strategiei alese.
  \item F3. Comutarea la runtime a backend-ului de cache, a strategiei, politicii, capacității și TTL-ului.
  \item F4. Statistici: hit/miss, interogări DB, mărimea cozii write-behind, statistici per nod.
  \item F5. Documentație interactivă a API-ului (OpenAPI/Swagger) și o interfață de control.
\end{itemize}

\subsection{Cerințe nefuncționale}
\begin{itemize}
  \item N1. \textbf{Performanță}: latență p50 sub 1 ms pentru hit-uri; debit de ordinul zecilor de mii de
        cereri/s pe un singur host.
  \item N2. \textbf{Scalabilitate}: adăugarea de instanțe web și de noduri de cache crește capacitatea.
  \item N3. \textbf{Coerență}: o scriere confirmată trebuie să fie vizibilă pe toate instanțele în cel mult
        câteva sute de milisecunde (strategiile sincrone) sau după flush (write-behind).
  \item N4. \textbf{Reziliență}: căderea cache-ului nu trebuie să oprească aplicația (degradare la DB);
        oprirea aplicației nu trebuie să piardă scrierile din coada write-behind.
  \item N5. \textbf{Reproductibilitate}: întregul mediu pornește cu o singură comandă; experimentele sunt
        scriptate.
\end{itemize}

\section{Arhitectura secvențială}
\label{sec:secvential}

Varianta de referință (\cref{fig:seqpar}a) este un singur proces care tratează cererile una câte una și
interoghează baza de date pentru fiecare citire. Analiza ei arată limitele:

\begin{itemize}
  \item debitul este $\lambda = 1/W$, unde $W$ include rețeaua și execuția interogării; cu $W \approx 1$ ms
        obținem cel mult $\approx 1.000$ cereri/s, indiferent de numărul de nuclee disponibile;
  \item baza de date este solicitată de fiecare cerere, deși 80–90\% dintre cereri vizează aceleași câteva
        mii de obiecte;
  \item nu există toleranță la defecte: căderea procesului oprește serviciul.
\end{itemize}

Varianta secvențială este reprodusă experimental prin backend-ul \texttt{none} cu un singur client
(experimentul E4, concurență 1) și, pentru motorul de cache, prin configurația 1 fir × 1 shard.

\section{Arhitectura paralelă și distribuită}
\label{sec:paralel}

\begin{figure}[H]
  \centering
  \resizebox{0.95\linewidth}{!}{% Arhitectura secvențială vs paralelă/distribuită.
\begin{tikzpicture}[node distance=5mm and 7mm]
  \node[font=\small\bfseries] (t1) {(a) Secvențial};
  \node[comp, below=3mm of t1] (c1) {1 client};
  \node[comp, below=of c1, minimum width=28mm] (p1) {1 proces, 1 fir\\[-1pt]\scriptsize cerere după cerere};
  \node[store, below=of p1] (d1) {PostgreSQL};
  \draw[arr] (c1) -- (p1); \draw[arr] (p1) -- node[note, right] {fiecare citire} (d1);

  \node[font=\small\bfseries, right=50mm of t1] (t2) {(b) Paralel și distribuit};
  \node[comp, below=3mm of t2, minimum width=36mm] (c2) {N clienți concurenți};
  \node[comp, below=of c2, minimum width=36mm, fill=c1!12] (lb) {load balancer};
  \node[comp, below left=5mm and -12mm of lb, minimum width=24mm, fill=c3!10] (i1) {instanță 1\\[-1pt]\scriptsize M task-uri tokio\\[-2pt]\scriptsize K fire};
  \node[comp, below right=5mm and -12mm of lb, minimum width=24mm, fill=c3!10] (i2) {instanță 2\\[-1pt]\scriptsize M task-uri tokio\\[-2pt]\scriptsize K fire};
  \node[comp, below=12mm of $(i1)!0.5!(i2)$, minimum width=42mm, fill=c2!12] (cc) {cache distribuit (S shard-uri × noduri)};
  \node[store, below=of cc] (d2) {PostgreSQL};
  \draw[arr] (c2) -- (lb); \draw[arr] (lb) -- (i1); \draw[arr] (lb) -- (i2);
  \draw[arr] (i1) -- (cc); \draw[arr] (i2) -- (cc);
  \draw[arr] (cc) -- node[note, right] {doar la miss} (d2);
\end{tikzpicture}
}
  \caption{Arhitectura secvențială (a) față de arhitectura paralelă și distribuită (b).}
  \label{fig:seqpar}
\end{figure}

Arhitectura propusă (\cref{fig:seqpar}b) introduce paralelism pe mai multe niveluri:

\begin{table}[H]
\centering\small
\begin{tabularx}{\linewidth}{lXl}
\toprule
Nivel & Mecanism & Tip \\
\midrule
Cereri & runtime asincron tokio: mii de task-uri multiplexate pe un pool de fire (work-stealing) & concurență + paralelism \\
Structura de date & \texttt{memtable}: $S$ shard-uri independente, fiecare cu propriul mutex & paralelism pe date \\
Instanțe & 2 instanțe web identice, fără stare, în spatele nginx (\texttt{least\_conn}) & scalare orizontală \\
Cache distribuit & $n$ noduri, chei partiționate prin consistent hashing & partiționare (sharding) \\
Invalidare & difuzare paralelă a invalidărilor pe tag (\texttt{JoinSet}) și \texttt{NOTIFY} & comunicare distribuită \\
Scrieri & write-behind: 4 workeri asincroni, partiționați după cheie, grupează scrierile în loturi & pipeline producător–consumator \\
\bottomrule
\end{tabularx}
\caption{Nivelurile de paralelism și distribuire.}
\end{table}

\subsection{Descompunerea problemei}

Cererile HTTP sunt independente (\emph{embarrassingly parallel}); singurele puncte de sincronizare sunt
structurile partajate: cache-ul și baza de date. Sharding-ul elimină sincronizarea globală în cache:
două chei din shard-uri diferite pot fi accesate simultan fără coordonare. Pentru cache-ul distribuit,
partiționarea prin consistent hashing face ca fiecare cheie să aibă un singur proprietar, deci nu este
nevoie de protocoale de consens; singura operație care atinge toate nodurile este invalidarea pe tag.

\subsection{Comunicare și sincronizare}

\begin{itemize}
  \item instanță ↔ nod de cache: protocol text binary-safe peste TCP, conexiuni persistente din pool,
        suport pentru \emph{pipelining};
  \item instanță ↔ PostgreSQL: pool de conexiuni (sqlx); notificări asincrone prin \texttt{LISTEN};
  \item în proces: \texttt{Mutex} per shard, \texttt{RwLock} pentru configurația curentă (înlocuită
        atomic prin \texttt{Arc}), canal \texttt{mpsc} mărginit pentru write-behind (\emph{backpressure}),
        \texttt{OnceCell} per cheie pentru single-flight.
\end{itemize}

\subsection{Compromisuri}

Distribuirea aduce costuri noi: o ședere în rețea per acces (≈100–300 µs față de ≈1 µs pentru L1),
problema coerenței între instanțe și complexitate operațională. De aceea sistemul oferă și varianta pe
două niveluri (L1 local cu TTL scurt + L2 distribuit), care combină latența cache-ului local cu
capacitatea agregată a clusterului.
